% sections/appendix_findings.tex
\section{Campaign findings}
\label{app:findings}

This appendix holds the evidence behind C3 (\cref{sec:findings}) and the GPU proxy of C2 (\cref{sec:sim}). It
breaks the campaign's step-count gains down by kind (\cref{sec:steps}), explains how to read
\cref{tab:levers} (\cref{sec:lever-notes}), examines the row pool, attention-source reuse and the EMA prewarm
(\cref{sec:rowpool,sec:asr,sec:prewarm}) and the noise and selection effects (\cref{sec:seeds,sec:salts}), and
lists what did not work (\cref{sec:negative}). It then asks how constant the price of a step is (\cref{app:kappa}),
prices the remaining gap (\cref{sec:gap}), describes the GPU proxy (\cref{sec:gpu}) and closes with further related
work (\cref{app:related-recipe}).

\subsection{Where the steps came from}
\label{sec:steps}

On 18~September a 30-minute run completed \val{steps-ktwentyfour} optimizer steps (\upl{K24}, depth~12, chip~A);
the final recipe completes \val{best-steps} (\upl{K82s4}, depth~9, chip~H). That ratio crosses chips, depths and batch
schedules, so it is not one effect. \Cref{tab:steps} breaks it into comparisons between uploads rehearsed on the
same chip. Three kinds of change contributed. (i) \emph{Host dispatch and kernel work at a fixed model and batch},
the only kind that would carry over to any recipe: on chip~B, \upl{K43}$\to$\upl{K46} added \val{sd-k43-k46} steps,
and the same-recipe step-time wins we measured were a scalar stack (\val{win-scalar}\docnum{}), a fused cross-entropy
(\val{win-fusedce}\docnum{}), sharding Muon (\val{win-muonshard}\docnum{}) and fusing the optimizer update (about
\val{win-optfuse}\docnum{}). (ii) \emph{A smaller model}: depth 10$\to$9 added \val{sd-k40-k41} on chip~A.
(iii) \emph{A batch schedule} that runs 65k- and 131k-token steps early, which adds optimizer steps that are not
full-batch steps by construction (\val{sd-k46-k50} on chip~B, together with other changes). The large early jump on
chip~A mixes FLOP reductions (grouped-query attention) with dispatch work. \Cref{fig:steps} (\cref{app:ledger}) shows
every upload's step count.

\begin{table}[!htbp]
  \centering
  \caption{\textbf{Step-count changes between uploads rehearsed on the same chip} (rehearsals, cold except
    \upl{K38}'s, \texttt{results/official-scores.csv}). Each row bundles every change between the two uploads.}
  \label{tab:steps}
  % GENERATED by paper/analysis -- steps between same-chip uploads (oracle.py)
{\footnotesize
\begin{tabular}{@{}llrrlp{5.6cm}@{}}
\toprule
Chip & Uploads & Steps & Change & Kind & What changed \\
\midrule
A & \upl{K24} $\to$ \upl{K38} & 912 $\to$ 1{,}526 & \ensuremath{+}67.3\% & FLOPs and dispatch & stream rows, value embeddings, GQA (one KV head), head size 256, scalar stack \\
A & \upl{K40} $\to$ \upl{K41} & 1{,}393 $\to$ 1{,}507 & \ensuremath{+}8.2\% & model size & depth 10 $\to$ 9 \\
B & \upl{K43} $\to$ \upl{K46} & 1{,}610 $\to$ 1{,}753 & \ensuremath{+}8.9\% & dispatch & scalar constants, async optimizer, Muon views, fused update \\
B & \upl{K46} $\to$ \upl{K50} & 1{,}753 $\to$ 1{,}943 & \ensuremath{+}10.8\% & batch schedule & batch warm-up 131k$\to$262k (+ MTP, Muon schedule) \\
C & \upl{K51} $\to$ \upl{K56} & 2{,}036 $\to$ 2{,}360 & \ensuremath{+}15.9\% & dispatch and schedule & ZeRO Muon/AdamW, 3-phase batch ramp, accumulation in the graph \\
C & \upl{K56} $\to$ \upl{K63} & 2{,}360 $\to$ 2{,}385 & \ensuremath{+}1.1\% & budget & time target 1,795 s, EMA off, cooldown 0.7 \\
\bottomrule
\end{tabular}
}

\end{table}

\subsection{Reading the levers table}
\label{sec:lever-notes}

\paragraph{Survivors and their noise scale.} \Cref{tab:levers} lists survivors. Its levers were kept because their
measured difference was favourable, out of many changes tried (\cref{tab:negative}; none of
\val{scalar-keys}\docnum{} scalar knob keys won), so each measured $\Delta$ is biased away from zero by the same
winner's curse as a salt lottery (\cref{sec:salts}), most for single pairs near the noise scale, that is, for the
late levers. The noise scale is itself uncertain: it rests on a salt SD quoted as a range (\val{salt-sd}\docnum{}),
and a change that acts from step~0, as every architecture lever does, may re-roll the warm-up path and then faces the
seed SD, \val{seed-sd-all-lo}--\val{seed-sd-all-hi} per run including spike paths (\cref{sec:seeds}), up to several
times the salt SD. Against that scale even attention-source reuse is at most a few SDs from zero, so the ratios in
the table are upper bounds for the step-0 levers.

\paragraph{Three entries.}

Three entries of \cref{tab:levers} need a word. \emph{LeakyReLU(0.35)$^2$} appears with three sizes because they are
three comparisons: \val{lever-leaky}\docnum{} in the campaign log on an earlier base, \val{leaky-chip-c}\docnum{} in a
chip-C pair against ReLU$^2$ at equal steps (the comparison the GPU proxy reproduces), and \val{leaky-official}
officially (\upl{K59}$-$\upl{K57}). \emph{Cooldown~0.7} is reported in the table together with the lower small-batch learning
rate, as three step-normalised chip pairs\docnum{}, all three negative; \cref{tab:gpu} uses two raw chip pairs of the
cooldown change alone, of opposite sign. The official \upl{K63}$-$\upl{K60} delta (\val{cooldown-official}) is
confounded, because \upl{K63} also turned the EMA off and added 2~s of budget. \emph{Official deltas} between uploads
published to four decimals carry a rounding error of up to $\pm2\times$\val{official-rounding}.

\subsection{The row pool: decorrelating a streaming loader}
\label{sec:rowpool}

\paragraph{Mechanism.} The kit's loader fills each row of a micro-batch from one of several parallel token streams,
so the rows of consecutive micro-batches are contiguous slices of the same few streams. From
micro-batch~\val{rowpool-from} on, our loader instead draws rows without replacement from a look-ahead pool of
\val{rowpool-size} micro-batches (\knob{FF_ROW_SHUFFLE}), a shuffle buffer at row granularity
\citep[cf.][]{xu2022corgipile,mishchenko2020reshuffling}. The first \val{rowpool-from} micro-batches, which cover the
warm-up, are unchanged, so that the seed's warm-up path (\cref{sec:seeds}) is not re-rolled. The shuffle has its own
RNG with a salt, which re-draws the order after micro-batch~\val{rowpool-from} without touching initialisation.

\paragraph{Evidence.} On chip E, the pool scored \val{rowpool-chip-treat}\docnum{} at \val{rowpool-steps-treat}\docnum{}
steps against \val{rowpool-chip-ctrl}\docnum{} at \val{rowpool-steps-ctrl}\docnum{} for the same file without it:
\val{rowpool-chip}\docnum{} raw, about \val{rowpool-stepnorm}\docnum{} after normalising for the
\val{rowpool-step-cost}\docnum{} of steps the pool cost, and about \val{rowpool-z} noise SDs. Officially, \upl{K70a}
scored \val{official-k70a} against \upl{K63}'s \val{official-k63}, a difference of \val{rowpool-official}. That
agreement is what determinism predicts and is not independent evidence that the effect generalises: both numbers come
from the same trajectory pair at seed~73. And because the pool changes the data order, its measured effect contains
one draw of the order noise (SD about \val{salt-sd}\docnum{}). A pool of \val{pool-large}\docnum{} gave the same
result as \val{rowpool-size}; a pool of \val{pool-small}\docnum{} helped less. Other order manipulations were null
within $\pm\val{rowpool-null}$\docnum{}: larger pools, epoch reshuffles, pool-aware accumulation and document
reordering.

\paragraph{What we do not know.} The evidence is one chip pair at one seed and one salt. It is probably real at about
\val{rowpool-z} order-noise SDs, but it is not replicated, and the pair is not in the released run data. Our working hypothesis is
that the pool reduces within-batch correlation between rows of the same stream, and with it gradient variance, in the
spirit of shuffling results for SGD \citep{gurbuzbalaban2019reshuffling}. We have \emph{not} measured this; the checks
are listed in \cref{sec:limits}.

\subsection{Attention-source reuse}
\label{sec:asr}

Blocks 6, 7 and 8 reuse the attention source computed by block~5 (\knob{FF_ATTN_SRC=5:6,7,8}), in the spirit of
cross-layer attention sharing \citep{brandon2024cla,xiao2019sharing}. It is our best-replicated structural effect:
four same-seed chip pairs (seeds 73, 67 and 58 on chip~C, seed~73 on chip~D) gave \val{asr-pairs}\docnum{} (mean
\val{asr-mean}), at a step cost under \val{asr-step-cost}\docnum{}; a fifth pair, seed~97 on chip~D, was
\val{asr-s97-delta}\docnum{} raw, mostly a \val{asr-s97-step-deficit}\docnum{}-step deficit (\cref{tab:postfit}).
Officially, \upl{K60} improved on \upl{K59} by \val{asr-official}. The surrogate, whose coefficient for this feature
then rested on two runs, predicted the right sign on four of five post-fit attention-source pairs but overestimated
the size by about \val{asr-overestimate}\docnum{} (\cref{tab:postfit}).

\subsection{The EMA blend and the prewarm accounting trap}
\label{sec:prewarm}

\paragraph{The trap.} Our final weights are a 0.6 blend of the live weights with an exponential moving average
\citep{polyak1992averaging,moralesbrotons2024ema}. The EMA update is a separate graph. On a fresh instance, the first
EMA update compiles that graph \emph{inside} the charged clock and costs \val{ema-compile-steps}\docnum{} steps. A warm
rehearsal, run after an earlier EMA run on the same instance, finds the graph in the compiler cache and does not pay.
Warm rehearsals therefore looked \val{ema-warm-bias}\docnum{} bpb better than the cold official run would be.

\paragraph{A regression, not a discovery.} The fix, \knob{FF_EMA_PREWARM}, runs one EMA update during the excluded
start-up. It was not new on the final day: the ledger lists the EMA prewarm already in \upl{\val{first-prewarm}}
(26~September). It disappeared when \upl{\val{prewarm-dropped}} turned the EMA off, and \upl{K77a} re-enabled the EMA
without it. Only cold rehearsals exposed the regression, which is the general lesson below.

\paragraph{Evidence.} Step counts compared within one instance show the cost (\cref{tab:prewarm}). On chip~G the recipe without EMA
ran \val{steps-k73s6} steps, with EMA but without prewarm \val{steps-k77a}, and with both \val{steps-k82s7}; on chip~H,
without EMA \val{steps-k73s4} and with both \val{steps-k82s4}. These are single runs, and step counts jitter by
\val{jitter-range-lo}--\val{jitter-range-hi} steps within a seed sweep (SD \val{jitter-sd-lo}--\val{jitter-sd-hi};
\cref{sec:determinism}). The chip-G drop of \val{prewarm-g-drop} steps without prewarm is about \val{prewarm-g-drop-sds}
single-run jitter SDs but only just above the largest within-sweep range (\val{jitter-range-hi} steps), while the
difference between prewarm and no EMA (\val{prewarm-g-gain} steps on chip G, \val{prewarm-h-gain} on chip H) is inside
it. The salts differ between some of these files, which changes the data order but not the step time.
Officially, \upl{K82s4} (prewarm) scored \val{prewarm-official} against \upl{K77a} (no prewarm): one pair, below the
oracle's resolution (\cref{sec:resolution}), consistent with the step-based estimate. The EMA blend itself is a separate
effect of about \val{ema-blend-chip}\docnum{} in one chip pair without prewarm; officially \upl{K77a} improved on
\upl{K73s4} by \val{ema-blend-official}.

\begin{table}[!htbp]
  \centering
  \caption{\textbf{Step counts around the compile trap} (single cold rehearsals, from
    \texttt{results/official-scores.csv}). Compare within a chip only; within-chip jitter is
    \val{jitter-range-lo}--\val{jitter-range-hi} steps.}
  \label{tab:prewarm}
  \small
  \begin{tabular}{llccr}
    \toprule
    Chip & Upload & EMA blend & EMA prewarm & Steps \\
    \midrule
    \val{chip-k73s6} & \upl{K73s6} & -- & -- & \val{steps-k73s6} \\
    \val{chip-k77a} & \upl{K77a} & \checkmark & -- & \val{steps-k77a} \\
    \val{chip-k82s7} & \upl{K82s7} & \checkmark & \checkmark & \val{steps-k82s7} \\
    \midrule
    \val{chip-k73s4} & \upl{K73s4} & -- & -- & \val{steps-k73s4} \\
    \val{chip-k82s4} & \upl{K82s4} & \checkmark & \checkmark & \val{steps-k82s4} \\
    \bottomrule
  \end{tabular}
\end{table}

\paragraph{The general lesson.} Any lazily compiled path that first runs after the clock starts is a hidden cost: EMA
updates, evaluation-mode switches, a late schedule phase with a new batch shape. On a compiled accelerator the compile
cost can exceed the benefit of the feature that triggers it, and only a \emph{cold} rehearsal reveals it. We recommend
that speed-run harnesses either compile all graphs before the clock starts or report compile time separately.

\subsection{Seeds and the warm-up coin flip}
\label{sec:seeds}

The largest single source of variance in our data was not a setting. In warm-up steps 3--10, some seeds take a
\emph{spike path}: the loss falls faster to step~4 and then much more slowly, ending step~10 roughly
\val{spike-gap-nats} nats higher than the good path (\cref{fig:warmup}a). The logged training loss is the weighted
multi-token-prediction sum (weights 1/0.5/0.25), not per-token cross-entropy, which is why it is near 13~nats. We
classify a run as spike-path when its step-10 loss is at least \val{spike-l10} nats. Good-path runs of these recipes
reach \val{l10-good} nats at step~10 and spike-path runs \val{l10-spike}; we chose the threshold after seeing the data,
in the otherwise empty gap between them, and one run at \val{l10-mid} is reported as intermediate. On the three clean seed sweeps
of a fixed recipe in our data (\cref{tab:seeds}), \val{spike-k} of \val{spike-n} seeds with a logged curve took the
spike path, \val{spike-frac} (Wilson 95\% interval \val{spike-wilson-lo}--\val{spike-wilson-hi}). The sweeps share
\val{spike-shared-seeds} seed, and the path may be a property of the seed, recipe and hardware together, so the runs
are not fully independent; counting each of the \val{spike-distinct-seeds} distinct seeds once gives
\val{spike-seed-k} spike seeds (\val{spike-seed-wilson}). Spike-path runs ended \val{spike-delta} worse on average
(95\% CI \val{spike-delta-lo} to \val{spike-delta-hi}), smaller than the \val{spike-penalty-doc}\docnum{} we had quoted
during the campaign from a broader, less controlled set of about \val{harvested-runs-doc}\docnum{} runs. We consider the
controlled estimate the better one.

\begin{figure}[!htbp]
  \centering
  \includegraphics{warmup_paths.pdf}
  \caption{\textbf{The warm-up coin flip.} (a) Training loss (the weighted multi-token sum) over the first 20 steps
    for the \val{g7-nseeds} seeds of one recipe on one chip: spike-path seeds (dashed, triangles) drop faster to
    step~4, then stall. (b) Final score by path in the three seed sweeps; hollow points have no logged curve.}
  \label{fig:warmup}
\end{figure}

\paragraph{The variance budget.} Seeds re-draw the initialisation, the warm-up path and the data order. Within a
sweep, the seed SD is \val{seed-sd-all-lo}--\val{seed-sd-all-hi} including spike paths and
\val{seed-sd-good-lo}--\val{seed-sd-good-hi} among good-path seeds only. Shuffle salts re-draw only the order after
micro-batch~\val{rowpool-from}, with an SD of about \val{salt-sd}\docnum{}. Both are as large as most late
improvements (\cref{tab:levers}).

\paragraph{Selection.} Our seed, 73, was itself chosen by a lottery. In the 13-seed sweep of the \upl{K57} recipe it
ranks \val{g7-s73-rank} of \val{g7-nseeds}, \val{g7-s73-gap} better than the good-path mean, and every later
same-seed comparison is measured on that lottery winner. Effects measured on a selected seed need not hold at other
seeds \citep{cawley2010overfitting}. We followed two rules: a lever must act after step~19, so that it cannot re-roll
the warm-up path, and an effect under 0.0003 needs three seeds; the second was often not affordable. Our one official
test of a change to the first steps illustrates the first rule: an attention-input shift in blocks 6--8 (\upl{K65a})
scored \val{xshift-official} against its base. It may have re-rolled the warm-up path, but we did not log its
step-10 loss, and \val{xshift-official} is about \val{xshift-spike-ratio} times the controlled spike-path penalty
(\val{spike-delta}), so a genuinely harmful change cannot be ruled out.

\subsection{Salt lotteries and the winner's curse}
\label{sec:salts}

Because a rehearsal predicts its official score up to an offset, a lottery over shuffle salts can be run on chip
and the best rehearsal uploaded.
Four salts of the \upl{K73} recipe spanned \val{salt-span}\docnum{} on chip, and salt~4 was uploaded as \upl{K73s4};
the later uploads \upl{K77a} and \upl{K82s4} inherit it. A second lottery on the final recipe picked salt~7
(\upl{K82s7}); its size is not recorded in the released files. Where a lottery's draws ran on different chips, salt
and chip speed are confounded. Selection was only partly rewarded. Salt~4 kept its edge officially. Salt~6
rehearsed \val{salt6-reh-vs-salt4} against it, but on chip~G against chip~H, so that edge is confounded with the chip;
officially it scored \val{salt6-vs-salt4} against salt~4 (\upl{K73s6}). Salt~7 gave the best rehearsal of the campaign
(\val{salt7-rehearsal}) and tied salt~4 officially (\val{salt7-vs-salt4}).

A simple model is consistent with this pattern. A salt's rehearsal advantage is its true effect plus a part that
does not transfer to the official run, with SD $s_\eta$. Noise that enters only the official run widens the outcome
but does not bias a selected draw, so $s_\eta$ holds only what is specific to the rehearsal: the variation of the
text gap between models (SD \val{gap-sd}, \cref{sec:contest}) and the rehearsal's step-count jitter priced at
$\kappa$ (\val{wc-jit-lo}--\val{wc-jit-hi}), together \val{wc-eta-lo}--\val{wc-eta-hi}. This assumes that the
rehearsal is otherwise repeatable: the unexplained part of the within-chip offset SD (\val{offdec-rest},
\cref{app:offset-detail}) is counted as official-side, and repeated cold runs of one file and seed would test that
assumption (\cref{app:derivations}). Against a spread of
rehearsal scores of \val{salt-sd}\docnum{}, the expected share of a selected advantage that survives, the reliability
ratio $\lambda$ (\cref{app:derivations}), is \val{wc-lam-lo}--\val{wc-lam-hi} when every draw ran on one chip. When
the draws ran on different chips, the between-chip offset SD (\val{wc-between}) adds to the part that does not
transfer, and $\lambda$ falls to \val{wc-lam-cross-lo}--\val{wc-lam-cross-hi}. Lotteries are worth running when draws
are cheap, but the projected gain should be shrunk by a share of about \val{wc-shrink-lo}--\val{wc-shrink-hi}
on one chip and \val{wc-shrink-cross-lo}--\val{wc-shrink-cross-hi} across chips (both over the range of the salt SD
and the jitter),
and if the official shard overlaps the public text (\cref{sec:contest}), part of what survives is selection on the
evaluation set itself \citep{dwork2015reusable,blum2015ladder}.

\subsection{What did not work}
\label{sec:negative}

\Cref{tab:negative} lists the changes that lost. Two patterns stand out. First, \emph{capacity changes lost in both
directions}: depth~10 and two key/value heads lost while running \val{depth10-steps}\docnum{} fewer steps, and depth~8
lost despite \val{depth8-steps}\docnum{} more steps. We read this as the 9-block model being a local optimum of the
capacity-for-steps trade at our kernels' speed, but it rests on one or two chip pairs without an equal-step comparison,
so we cannot say whether depth~10 lost \emph{only} because of its steps. Second, \emph{scalar retuning found
nothing}: none of \val{scalar-keys}\docnum{} scalar knob keys (learning-rate scales, betas, weight decay,
initialisation, momentum, soft cap, RoPE base, cooldown $\pm0.1$) won a confirmation. A single chip pair's noise SD
(\val{noise-chip-pair}) is as large as the effect a retune could plausibly have, so unless a key was tried several
times this is an absence of evidence for scalar gains, not evidence that none exist. The four official tests on the final day
(\upl{K65a}--\upl{K65d}) were all worse than or tied with their base.

\begin{table}[!htbp]
  \centering
  \caption{\textbf{What lost.} $\Delta$ against the control (positive is worse); evidence as in \cref{tab:levers}.
    Official deltas are against \upl{K63}. \textsuperscript{\S}\,May have re-rolled the warm-up path; its step-10 loss
    was not logged (\cref{sec:seeds}). Source: \texttt{docs/FINDINGS.md} \S5\docnum{} and
    \texttt{results/official-scores.csv}.}
  \label{tab:negative}
  % GENERATED by paper/analysis -- what lost, docs/FINDINGS.md section 5 and results/official-scores.csv (K65a-d)
{\small
\begin{tabular}{@{}p{7.6cm}rlr@{}}
\toprule
Change & $\Delta$ ($10^{-3}$) & Evidence & $n$ \\
\midrule
Depth 10 (bigger; 6--11\% fewer steps) & \ensuremath{+}3.8 to \ensuremath{+}4.3 & chip pair & 2 \\
Two KV heads (bigger) & \ensuremath{+}5.7 & chip pair & 1 \\
Depth 8 (smaller; +10.6\% steps) & \ensuremath{+}3.0 to \ensuremath{+}5.0 & chip pair & 2 \\
Late $k{=}8$ batch in the cooldown (K65b) & \ensuremath{+}0.5 & official pair & 1 \\
Split AdamW cooldown 0.4 (K65d) & \ensuremath{+}0.6 & official pair & 1 \\
EMA-Nesterov lookahead 0.3 (K65c) & \ensuremath{+}0.04 & official pair & 1 \\
AdamW LR $-$16\% & \ensuremath{+}0.8 to \ensuremath{+}0.9 & chip pair & 2 \\
AdamW warm-up reshaping (steps 0--19) & \ensuremath{+}1.0 & chip pair & 1 \\
Drop the block-0 MLP & \ensuremath{+}1.9 & chip pair & 1 \\
No positional encoding in layers 3--4 & \ensuremath{+}3.3 & chip pair & 1 \\
Attention-input shift, blocks 6--8 (K65a)$^{\S}$ & \ensuremath{+}3.4 & official pair & 1 \\
Scalar retunes (LR scales, betas, WD, init, momentum, soft cap, RoPE base, cooldown $\pm$0.1) & 0 of 61 keys won & chip pairs & 61 \\
Bigger pools, epoch reshuffles, pool-aware accumulation, document reordering & 0 $\pm$ 0.3 & chip pairs & -- \\
\bottomrule
\end{tabular}
}

\end{table}

\subsection{How constant the price of a step is}
\label{app:kappa}

The law of \cref{eq:exchange} is not constant across the campaign. The simulator's prior
(\cref{sec:quality}) gives it a curvature of $c=\val{q-prior-c}$ per $(\ln S)^2$, so that the local slope is
$\kappa(S)=\val{kappa}-2c\ln(S/2300)$, which is \val{kappa-at-anchor} at the geometric mid-point of the long-run calibration if its shorter run ran at
\upl{K44}'s step count, and the campaign log records \val{slope-900}\docnum{} at about 900 steps. That curvature is
itself a prior set from these campaign rates, so it is consistent with the two calibrations rather than evidence for
them. The surrogate's fitted slope (\val{q-slope-pct} per 1\%, \cref{sec:quality}) is \emph{not} independent
evidence either: its prior mean is $\kappa$ itself. Nor do the run data pin $\kappa$ down: refitting the surrogate with
a vague slope prior (SD \val{kappa-vague-prior-sd} instead of \val{q-prior-b-sd}) gives $\kappa=\val{kappa-vague}$
(95\% interval \val{kappa-vague-lo}--\val{kappa-vague-hi}), which excludes the value we use. That fit's slope is
confounded with mechanisms that change both steps and quality, so it is not a better estimate, but at its interval
the gap to \#10 would cost \val{gap-ten-vague-range} more compute instead of \val{gap-ten-range}, and a typical
final-day change of \val{late-change}\docnum{}~bpb would be worth \val{late-change-vague-range} more steps instead of
\val{late-change-pct}.

\subsection{Pricing the gap in compute}
\label{sec:gap}

\Cref{tab:gap} converts the distance from our best official score, \val{best-official}, to reference points on the final
public leaderboard snapshot (about \val{leaderboard-time}\docnum{}) into effective compute through
\cref{eq:exchange}. Closing the \val{gap-ten} gap to \#10 would have needed about \val{gap-ten-range} more effective
compute across our two calibrations of $\kappa$. The gap to \#1, \val{gap-one-range}, extrapolates far beyond the
\val{xr-anchor-pct} anchor and is an order of magnitude only. \Cref{fig:exchange} draws the price of a step with both
calibrations and these gaps.

\begin{table}[!htbp]
  \centering
  \caption{\textbf{What each leaderboard reference would cost in effective compute}, from our best official score
    (\val{best-official}, \upl{K82s4}), as $\exp(\text{gap}/\kappa)-1$ for the price used ($\kappa=\val{kappa}$) and
    for the long-run calibration. Leaderboard values are from the final public snapshot, referred to by rank.}
  \label{tab:gap}
  % GENERATED by paper/analysis -- compute needed = exp(gap / kappa) - 1 (exchange.py)
{\small
\begin{tabular}{@{}lrrrr@{}}
\toprule
Target & val\_bpb & Gap & Compute, $\kappa=0.057$ & Compute, $\kappa=0.063$ \\
\midrule
\#10 on the final snapshot & 0.9554 & 0.0060 & +11\% & +10\% \\
\#1 on the final snapshot & 0.9328 & 0.0286 & +65\% & +58\% \\
\bottomrule
\end{tabular}
}

\end{table}

\begin{figure}[!htbp]
  \centering
  \includegraphics{exchange_rate.pdf}
  \caption{\textbf{The price of a step} (\cref{eq:exchange}). Solid line: $\kappa=\val{kappa}$, the price used
    throughout; dashed line: $\kappa=\val{xr-k-anchor}$ from the single long run (square). Black bars: the compute
    needed to close the gap from our best score to \#10 and \#1 on the final leaderboard snapshot, across the two
    values of $\kappa$. Triangle: a typical final-day recipe change of \val{late-change}\docnum{}~bpb, worth
    \val{late-change-pct} more steps.}
  \label{fig:exchange}
\end{figure}

\paragraph{A hypothesis, not a finding.} The final recipe trained at about \val{tokens-per-s}\docnum{} tokens/s on a
\val{params}\docnum{}-parameter model, which we estimated at roughly \val{mfu}\docnum{} model-FLOPs utilisation
\citep{chowdhery2023palm} with about \val{idle}\docnum{} of wall time idle on host dispatch; both are approximate
profiles. At an equal recipe, \cref{eq:exchange} values \val{tput-lo-pct}--\val{tput-hi-pct}\% more tokens per second
at \val{tput-range}~bpb, more than our whole gap to \#10, while late recipe changes were worth about
\val{late-change}\docnum{} each. We therefore \emph{hypothesise} that the cheapest way to close the gap was kernel
throughput: custom NKI kernels \citep{aws_nki_docs} for attention, the fused MLP and the optimizer updates. Nothing in
our data rules out the alternative that better-ranked teams had better recipes or data handling, and the conversion
assumes that every team sits on our compute curve. A test would be to write one kernel, measure its step time with a
short chip screen, price it with \cref{eq:exchange} and confirm it with a cold rehearsal; we did not reach that stage
in Phase~1.

\subsection{The GPU proxy}
\label{sec:gpu}

Besides the surrogate, \ffsim{} has a second route: it runs the actual contest \texttt{train.py} on a CUDA GPU along the chip's optimizer trajectory, step for step (\cref{app:gpu} describes its parts). It measures per-step quality only; step time must come from the chip.

\Cref{tab:gpu} compares it, run on an L40S GPU, with Trainium. Two of the four comparisons are informative equal-step
chip pairs (attention-source reuse and LeakyReLU$^2$), and the GPU matched them to within \val{gpu-maxdiff-inf}\docnum{}.
The cooldown-0.7 reference is the mean of two chip pairs of opposite sign, and the split-AdamW reference is an
official upload delta at a different step count, so neither is a test. The noise is larger than the agreement: in
default CUDA mode two runs of one configuration on two GPU boxes differed by \val{gpu-box-diff}\docnum{}, because the
flash-attention backward pass is nondeterministic, and the campaign documents do not record which of the four
comparisons ran in deterministic mode. Forcing deterministic algorithms made runs identical at every logged step, at
\val{gpu-tps-det}\docnum{} instead of \val{gpu-tps-default}\docnum{} tokens/s. Hardware can also flip the warm-up
path: seed~58 takes the good path on chip and a spike path on the GPU. With $n=\val{gpu-n-inf}$ informative pairs and
noise above the agreement, the proxy is a feasibility demonstration, consistent with the chip in sign and magnitude,
not a validated instrument. A run of the full schedule cost about \val{gpu-hours}\docnum{}~GPU-hours (about
\$\val{gpu-cost-run}\docnum{}).

\begin{table}[!htbp]
  \centering
  \caption{\textbf{GPU proxy vs Trainium} at equal steps, seed 73 ($10^{-3}$ bpb except the last row; values from
    \texttt{docs/SIMULATOR.md}\docnum{}). \textsuperscript{*}\,Not an informative comparison: the Trainium value is
    the mean of two chip pairs of opposite sign (\val{gpu-cd-pairs}), or an official delta at a different step count.}
  \label{tab:gpu}
  % GENERATED by paper/analysis -- GPU proxy vs chip, 1e-3 bpb except the last row; values from docs/SIMULATOR.md route 2
{\footnotesize
\begin{tabular}{@{}p{4.6cm}rrrp{3.6cm}@{}}
\toprule
Comparison (seed 73, equal steps) & GPU & Trainium & Diff. & Trainium evidence \\
\midrule
Attention-source reuse (5:6,7,8) & \ensuremath{-}1.87 & \ensuremath{-}1.81 & \ensuremath{-}0.06 & chip C pair \\
LeakyReLU(0.35)$^2$ vs ReLU$^2$ & \ensuremath{-}1.14 & \ensuremath{-}1.05 & \ensuremath{-}0.09 & chip C pair \\
Cooldown 0.7$^{*}$ & \ensuremath{-}0.12 & \ensuremath{-}0.14 & \ensuremath{+}0.02 & mean of 2 chip pairs of opposite sign \\
Split AdamW cooldown 0.4$^{*}$ & \ensuremath{+}0.45 & \ensuremath{+}0.60 & \ensuremath{-}0.15 & official upload delta, other steps \\
\midrule
\upl{K60} absolute val\_bpb at 2{,}314 steps (bpb) & 0.959952 & 0.959948 & \ensuremath{+}0.004 & one run each \\
\bottomrule
\end{tabular}
}

\end{table}

\subsubsection{How the proxy is built}
\label{app:gpu}

The proxy is built from four parts. \emph{A generated fork}: a patcher applies \val{gpu-patches}\docnum{} checked
textual patches to the recipe, a check verifies that the committed fork is byte-identical to a fresh build, and the
patcher refuses any anchor it cannot find exactly. \emph{A virtual charged clock}: every clock-keyed switch
(accumulation phases, the eight multi-token-prediction stages, cooldown levels, the EMA save and the stop) fires at the
step number at which it fired on chip, computed from a chip log's per-phase step times; real GPU time is ignored.
\emph{Data-world emulation}: the trainer reproduces the organiser loader's eight-rank sharding on one device, so each
step sees the same tokens. \emph{Calibration before use}: no GPU verdict may influence a chip run until the proxy has
reproduced pairs the campaign already paid for; an earlier proxy was rejected under this rule when one of its verdicts
flipped sign. A GPU pair is void when its step-10 loss sits on a different warm-up path from its control. The shipped
\texttt{--deterministic} flag uses PyTorch's \texttt{warn\_only=True}; paired verdicts need \texttt{warn\_only=False}.
The protocol, runbook and patch notes are in \texttt{ffsim/gpu/}.

\subsection{Further related work}
\label{app:related-recipe}

\paragraph{Benchmarks, proxies and reproducibility.} DAWNBench \citep{coleman2017dawnbench} made time-to-accuracy an
end-to-end metric, and its analysis \citep{coleman2019dawnbench} found that metric to have a low coefficient of
variation and the speed-optimised models to generalise nearly as well as conventionally trained ones; the report on
the first AlgoPerf competition \citep{kasimbeg2025algoperf} discusses the engineering needed to compare submissions
fairly on fixed hardware. Both are organiser-side accounts of a benchmark. New ImageNet and CIFAR-10 test sets
\citep{recht2019imagenet} lowered accuracies while largely preserving the order of models, and out-of-distribution
accuracy tracks in-distribution accuracy closely across many distribution shifts \citep{miller2021accuracy}; the
rehearsal offset of \cref{sec:oracle}, a shift about the size of the 2M-to-20M text gap (\val{gap-mean}) that does
not depend on the recipe, is a small instance of the same pattern. For decisions in language-model development,
\citet{madaan2024variance} quantify the seed variance of evaluation benchmarks, \citet{heineman2025signal} relate a
benchmark's usefulness for small-scale decisions to its signal-to-noise ratio, and DataDecide
\citep{magnusson2025datadecide} measures how often small experiments pick the pretraining data that is best at a
larger scale. Surrogate benchmarks replace expensive evaluations with a regression model fitted to logged runs, for
hyperparameter optimisation \citep{eggensperger2015surrogates} and for architecture search
\citep{zela2022surrogate}. \citet{chen2022reproducible} make deep-learning training reproducible by recording and
replaying software randomness and patching hardware-related nondeterminism, and \citet{hoefler2015scientific} ask
authors to report whether measurements are deterministic, to give confidence intervals when they are not and to
compare them in a statistically sound way.

\paragraph{Optimizers, schedules and instability.} Our recipe uses Muon, momentum orthogonalised by a Newton--Schulz
iteration \citep{jordan2024muon,bjorck1971orthogonal}, for matrices and AdamW
\citep{kingma2015adam,loshchilov2019adamw} for the rest. Muon is analysed as steepest descent under a spectral norm
by \citet{bernstein2024anthology,bernstein2024modular}; scaling studies and variants include
\citet{liu2025muonscalable} and NorMuon \citep{li2025normuon}. Warm-up and decay follow the warmup-stable-decay and
cooldown literature \citep{loshchilov2017sgdr,hu2024minicpm,hagele2024scaling,defazio2024schedulefree}. Why warm-up
stabilises early training is analysed by \citet{liu2020radam} and \citet{xiong2020layernorm}; loss spikes and
instabilities in early or large-scale training are studied by \citet{wortsman2024proxies}, \citet{chowdhery2023palm},
\citet{molybog2023adam} and \citet{takase2025spike}. Our warm-up ``spike path'' (\cref{sec:seeds}) is a small-scale
instance. The batch warm-up follows the critical-batch-size view \citep{mccandlish2018largebatch} and
linear-scaling practice \citep{goyal2017imagenet}.

\paragraph{Weight averaging, data order and architecture.} Our final weights blend live weights with an
exponential moving average \citep{polyak1992averaging,izmailov2018swa,kaddour2022lawa,sanyal2023earlyavg,moralesbrotons2024ema,busbridge2023ema}.
Random reshuffling beats with-replacement sampling \citep{gurbuzbalaban2019reshuffling,mishchenko2020reshuffling},
partial shuffles trade I/O for statistical quality \citep{xu2022corgipile}, and packing and document boundaries
affect language-model quality \citep{krell2021packing,zhao2024sequence,ding2024truncations}; our row pool is a
row-level shuffle buffer on a streaming loader. The architecture uses attention \citep{vaswani2017attention} with
rotary embeddings \citep{su2021roformer}, grouped-query attention with one key/value head
\citep{shazeer2019mqa,ainslie2023gqa}, a squared-ReLU-family MLP \citep{so2021primer}, value embeddings related to
value-residual learning \citep{zhou2024valueresidual} and multi-token prediction \citep{gloeckle2024multitoken};
attention-source reuse is closest to cross-layer attention sharing \citep{brandon2024cla,xiao2019sharing}.
